\documentclass{subfiles}

\begin{document}
    \subsection{{\sArray}}
    {\sArray} adalah sebuah variabel yang menyimpan sekumpulan data yang memiliki tipe yang sama. Setiap data tersebut menempati lokasi atau alamat memori yang berbeda-beda dan selanjutnya disebut dengan elemen {\sarray}. Elemen {\sarray} itu kemudian dapat diakses melalui indeks yang terdapat di dalamnya. Namun, penting sekali untuk diperhatikan bahwa dalam C++ indeks {\sarray} selalu dimulai dari 0, bukan 1\parencite{raharjo2015pemrograman}.

    Dilihat dari dimensinya, {\sarray} dapat dibagi menjadi {\sarray} satu dimensi, {\sarray} dua dimensi dan {\sarray} multidimensi \parencite{faradiba2019penggunaan}:

    \subsubsection{{\sArray} satu dimensi}

    \begin{indentpar}\hpar{}
        Untuk mendeklarasikan variabel dengan tipe data {\sarray} satu dimensi pada notasi algoritma.
    \end{indentpar}
    \subsubsection{{\sArray} dua dimensi}

    \begin{indentpar}\hpar{}
        {\sArray} dua dimensi merupakan {\sarray} yang terdiri dari m buah baris (\textit{row}) dan n buah kolom (\textit{column}). Bentuk {\sarray} semacam ini menggunakan 2 (dua) buah kelompok indeks yang masing-masing direpresentasikan sebagai indeks baris dan kolom.
    \end{indentpar}
    \subsubsection{{\sArray} multidimensi}

    \begin{indentpar}\hpar{}
        Dalam menggambarkan {\sarray} multidimensi, hanya terbatas hingga dimensi ke-3, yakni dengan menggunakan bangun ruang, namun dalam kenyataannya, tipe data {\sarray} ini dapat dibentuk menjadi lebih dari tiga dimensi atau menjadi n-dimensi.
    \end{indentpar}

    \subsection{{\sString}}

    C++ memiliki dukungan terhadap dua jenis {\sstring}, yaitu {\sstring} gaya bahasa C dan gaya bahasa C++. {\sString} gaya bahasa C sering disebut dengan \textit{null terminated string} atau \textit{C-style string}. {\sString} jenis ini dapat diimplementasikan ke d dari alam dua bentuk. Bentuk pertama, {\sstring} diimplementasikan sebagai larik atau kumpulan dari tipe karakter (\textit{array of char}). Dalam cara ini, {\sstring} akan dianggap sebagai kumpulan karakter yang diakhiri oleh karakter \textit{null} (\citecode{\textbackslash 0}). Setiap karakter di dalam {\sstring} dapat diakses melalui indeks {\sarray}. Bentuk kedua, {\sstring} diimplementasikan sebagai \textit{pointer} ke tipe karakter (\textit{pointer  of char}). Berikut ini contoh kode yang menunjukan penggunaan {\sstring} dalam bahasa C \parencite{raharjo2015pemrograman}:

\begin{code}
\begin{verbatim}
char str1[100];    // C-string menggunakan array
char *str2;        // C-string menggunakan pointer
\end{verbatim}
\end{code}

    C++ mengimplementasikan string sebagai kelas. Kelas yang digunakan untuk merepresentasikan {\sstring} adalah \citecode{string} (huruf s ditulis dalam huruf kecil). Kelas tersebut berisi kumpulan data dan fungsi yang berguna untuk mengangani permasalahan-permasalahan program yang berkaitan dengan {\sstring} \parencite{raharjo2015pemrograman}.

    Fungsi string yang telah tersedia pada bahasa C yaitu sebagai berikut \parencite{sasongko2010modul}:
    \subsubsection{\textit{Strlen}}
    \begin{indentpar}\hpar{}
        Fungsi untuk mengetahui jumlah teks pada data karakter. Parameternya adalah variabel yang menandung teks, nilai baliknya adalah angka bernilai panjang teksnya.
    \end{indentpar}
    \subsubsection{\textit{Strcpy}}
    \begin{indentpar}
        Fungsi ini digunakan untuk menyalin nilai teks pada suatu variabel ke variabel lain.
    \end{indentpar}
    \subsubsection{\textit{Strcat}}
    \begin{indentpar}\hpar{}
        Fungsi ini digunakan untuk menggabungkan nilai \textit{text} suatu variabel maupun {\sstring} konstan dengan nilai teks lainnya.
    \end{indentpar}
\end{document}

